% =============================================================================
% OPENING
% =============================================================================

\begin{frame}[plain]
  \setbeamercolor{background canvas}{bg=navy}
  \begin{tikzpicture}[remember picture,overlay]
    \fill[navy] (current page.south west) rectangle (current page.north east);
    \fill[blue] ([xshift=-4.4cm]current page.north east) rectangle (current page.south east);
    \node[anchor=west,text=white,font=\bfseries\fontsize{34}{38}\selectfont] at ([xshift=0.9cm,yshift=1.0cm]current page.west) {Matrices};
    \node[anchor=west,text=white!78,font=\Large] at ([xshift=0.95cm,yshift=-0.15cm]current page.west) {Read the shape. Understand the action.};
    \node[anchor=west,text=white!65,font=\small] at ([xshift=0.95cm,yshift=-2.2cm]current page.west) {Basic Mathematics \;\textbullet\; Lecture 1};
    \node[text=white,font=\Large] at ([xshift=-2.15cm]current page.east) {$\begin{pmatrix}1&2&0\\-1&3&4\\2&0&5\end{pmatrix}$};
  \end{tikzpicture}
\end{frame}

\begin{frame}{The story of the lecture}
  \begin{center}
  \begin{tikzpicture}[node distance=0.55cm and 0.45cm]
    \tikzset{storybox/.style={rounded corners=5pt,minimum width=2.2cm,minimum height=1.0cm,align=center,font=\bfseries,draw=line,very thick}}
    \node[storybox,fill=softblue] (a) {Read\\the object};
    \node[storybox,fill=softgreen,right=of a] (b) {Operate\\entrywise};
    \node[storybox,fill=softorange,right=of b] (c) {Compose\\actions};
    \node[storybox,fill=softblue,right=of c] (d) {Notice\\structure};
    \node[storybox,fill=softgreen,right=of d] (e) {Transform\\the problem};
    \draw[-{Latex[length=3mm]},thick,blue] (a) -- (b);
    \draw[-{Latex[length=3mm]},thick,blue] (b) -- (c);
    \draw[-{Latex[length=3mm]},thick,blue] (c) -- (d);
    \draw[-{Latex[length=3mm]},thick,blue] (d) -- (e);
  \end{tikzpicture}
  \end{center}

  \vspace{0.7cm}
  \takeaway{A matrix is useful because it keeps structure visible while we calculate.}
\end{frame}

\begin{frame}{Three 45-minute arcs}
  \begin{columns}[T,onlytextwidth]
    \begin{column}{0.31\textwidth}
      \begin{block}{Arc I — Language}
        \centering
        shape, entries, rows, columns\\[0.3cm]
        addition and scaling
      \end{block}
    \end{column}
    \begin{column}{0.31\textwidth}
      \begin{exampleblock}{Arc II — Action}
        \centering
        row $\times$ column\\[0.3cm]
        composition and order
      \end{exampleblock}
    \end{column}
    \begin{column}{0.31\textwidth}
      \begin{alertblock}{Arc III — Structure}
        \centering
        identity, transpose, diagonal\\[0.3cm]
        row operations
      \end{alertblock}
    \end{column}
  \end{columns}
\end{frame}
